```latex
\begin{proof}
Fix a positive diagonal matrix
\[
D=\operatorname{diag}(d_1,\ldots,d_m),
\qquad d_i>0.
\]
Let
\[
A:=DHD^{-1}.
\]
We will prove that
\[
\|A\|_{r\to r}\ge m^{1-1/r}.
\]

Choose an index \(i^*\in\{1,\ldots,m\}\) such that
\[
d_{i^*}=\max_{1\le j\le m} d_j.
\]
Let \(h_{i^*}\in\{\pm 1\}^m\) denote the \(i^*\)-th row of \(H\), viewed as a column vector. By the stated property of \(H\), we have
\[
Hh_{i^*}=me_{i^*}.
\]
Now define the test vector
\[
y:=Dh_{i^*}.
\]
Since \(D\) is invertible and \(h_{i^*}\neq 0\), we have \(y\neq 0\), so it is a valid vector on which to test the operator norm.

We compute
\[
Ay
=
DHD^{-1}Dh_{i^*}
=
DHh_{i^*}
=
D(me_{i^*})
=
md_{i^*}e_{i^*}.
\]
Therefore
\[
\|Ay\|_r
=
\|md_{i^*}e_{i^*}\|_r
=
md_{i^*}.
\]

On the other hand, since every entry of \(h_{i^*}\) has absolute value \(1\), we have
\[
\|y\|_r
=
\|Dh_{i^*}\|_r
=
\left(\sum_{j=1}^m |d_j h_{i^*j}|^r\right)^{1/r}
=
\left(\sum_{j=1}^m d_j^r\right)^{1/r}.
\]
Because \(d_{i^*}=\max_j d_j\), we have \(d_j\le d_{i^*}\) for every \(j\). Hence
\[
\sum_{j=1}^m d_j^r
\le
\sum_{j=1}^m d_{i^*}^r
=
m d_{i^*}^r.
\]
Taking \(r\)-th roots gives
\[
\|y\|_r
=
\left(\sum_{j=1}^m d_j^r\right)^{1/r}
\le
m^{1/r}d_{i^*}.
\]

By the definition of the induced operator norm,
\[
\|A\|_{r\to r}
=
\sup_{x\neq 0}\frac{\|Ax\|_r}{\|x\|_r}.
\]
Since \(y\neq 0\), we may test the supremum at \(y\). Thus
\[
\|A\|_{r\to r}
\ge
\frac{\|Ay\|_r}{\|y\|_r}
=
\frac{md_{i^*}}{\left(\sum_{j=1}^m d_j^r\right)^{1/r}}
\ge
\frac{md_{i^*}}{m^{1/r}d_{i^*}}
=
m^{1-1/r}.
\]
Therefore
\[
\|DHD^{-1}\|_{r\to r}\ge m^{1-1/r}.
\]

Since the positive diagonal matrix \(D\) was arbitrary, every admissible diagonal scaling satisfies the same lower bound. Taking the infimum over all such \(D\), we obtain
\[
\rho_r(H)
=
\inf_{\substack{D=\operatorname{diag}(d_1,\ldots,d_m)\\ d_i>0}}
\|DHD^{-1}\|_{r\to r}
\ge
m^{1-1/r}.
\]
This proves the desired lower bound.
\end{proof}
```